\section{Related Work}
\label{sec:related}

Research on autism spectrum disorder (ASD) classification using resting-state functional magnetic resonance imaging (rs-fMRI) has gradually developed from statistical modelling and functional-connectivity representation. More recent research is focused on multisite learning, deep architectures, dynamic connectivity, multimodal fusion, and reliability-aware prediction. In early studies, methods are established for representing high-dimensional covariance and brain-network information. These methods are later extended through deep learning, multiple brain atlases, temporal modelling, and graph-based approaches to improve classification performance. Alongside these developments, calibration, interpretability, selective prediction, and statistical validation are increasingly used to improve the reliability of machine-learning predictions. Overall, the performance and reliability of ASD classification are shaped by several methodological choices, including how functional connectivity is represented, how variation across acquisition sites is controlled, and how prediction uncertainty is assessed.

\subsection{Foundations of Functional Connectome-Based Analysis}

The methodological basis for later connectome-based ASD classification is supported by several important statistical developments. Exact confidence limits for binomial outcomes are proposed in an early statistical study \citep{ClopperPearson1934}, providing a basis for the reliable estimation of classification error. Alongside the measurement of uncertainty, stable estimation of high-dimensional connectivity is also required. For this purpose, covariance estimation is improved through the Ledoit--Wolf shrinkage method \citep{LedoitWolf2004}. This method is particularly relevant to rs-fMRI, where covariance matrices are estimated from a limited number of time points across many brain regions. Reliable comparison between classification models is further supported by statistical procedures developed for classifier evaluation \citep{Demsar2006}. As multisite datasets are increasingly used, variation between acquisition sites also needs to be controlled. The ComBat empirical-Bayes framework is originally proposed for the correction of batch effects \citep{Johnsonetal2007}. About a decade later, the same framework is adapted to multi-site neuroimaging to reduce scanner and site-related variation \citep{Fortinetal2018}.


Building on these statistical foundations, the geometric structure of covariance matrices is considered more directly in later functional-connectivity research. Group-level covariance modelling is developed by treating covariance matrices according to their non-Euclidean geometry \citep{Varoquauxetal2010}. From this formulation, tangent-space representations are established for the analysis of functional connectivity and later become widely used in connectome-based prediction. A biological framework for interpreting these connectivity patterns is also provided through the identification of seven large-scale intrinsic brain networks \citep{Yeoetal2011}. Through this organisation, large numbers of individual connections \reviewhl{are interpretable} within broader functional systems.

These developments provide a strong basis for static functional-connectivity analysis, where connectivity is estimated across the full duration of a scan. The temporal organisation of these connections is then examined through dynamic functional-connectivity research, in which changes in network relationships over time are shown to provide additional information \citep{Hutchisonetal2013}. Functional connectivity is therefore increasingly represented through both spatial patterns across brain regions and temporal changes in these patterns.


\subsection{Multisite ASD Data and Reproducibility}

Large-scale ASD neuroimaging research is substantially advanced through the release of ABIDE \citep{DiMartinoetal2014}. Data from multiple international sites are combined, making larger ASD classification studies possible. This increased scale also introduces variation in scanner type, acquisition protocol, demographic composition, and site characteristics. As a result, prediction models \reviewhl{risk capturing} acquisition-related differences instead of ASD-related neural patterns.

The use of larger multisite datasets also requires high-dimensional brain information to be represented carefully. Regularised matrix regression is proposed for this purpose, where predictors are kept in matrix form so that relationships among features are preserved during modelling \citep{ZhouLi2014}. Careful analysis is also required when changes in connectivity over time are examined. Sliding-window estimates are shown to depend on window length, and apparent connectivity changes \reviewhl{are potentially produced} by analysis choices \citep{LeonardiVanDeVille2015}. Temporal information therefore needs to be evaluated with careful control of scan length and window design.

These methodological choices directly affect whether findings \reviewhl{are reproducible} across datasets and sites. Using multisite resting-state data, biomarker performance is shown to depend on preprocessing, feature construction, and validation design \citep{Abrahametal2017}. The reliability of time-varying connectivity is also examined through a review of dynamic functional-connectivity methods \citep{Pretietal2017}. Together, these studies show that stable ASD-related findings depend on careful control of participants, analysis choices, and acquisition sites.

\subsection{Reliability, Calibration, and Explainability}

Alongside neuroimaging research on representation and generalisation, reliability-oriented machine-learning methods are also developed. Probability calibration is examined in a study \citep{Guoetal2017}, where high classification accuracy is shown to be insufficient for reliable confidence estimates. A model \reviewhl{potentially ranks} participants correctly while producing probabilities that do not match actual prediction frequencies. Selective classification is introduced in a study \citep{GeifmanElYaniv2017}, where uncertain cases are allowed to be rejected rather than requiring a prediction for every input. A trade-off between coverage and prediction error is therefore created, which is particularly relevant to screening applications where uncertain cases \reviewhl{are referred} for further clinical assessment. Interpretability is further supported through SHAP \citep{LundbergLee2017} to measure the contribution of individual features to each prediction. Evaluation is therefore extended beyond conventional accuracy, as calibration is used to assess prediction confidence, selective prediction is used to determine \reviewhl{the conditions under which a prediction is issued}, and explanation methods are used to identify the contributing features. Clinical usefulness is also examined through decision-curve analysis \citep{VanCalsteretal2018}, where statistical performance and practical decision value are evaluated separately. Reliable ASD classification is therefore increasingly evaluated through prediction accuracy, confidence, selective decision-making, interpretability, and clinical usefulness.


\subsection{Deep Learning for ASD Classification}

Alongside these reliability-focused developments, deep learning is increasingly applied to ABIDE functional-connectivity data. Distributed connectivity patterns are learned automatically and an accuracy of approximately 70.0\% is reported, showing that ASD-related information \reviewhl{is captured} through deep feature learning \citep{Heinsfeldetal2018}. The effect of multisite variation still needs to be addressed when such models are applied to ABIDE. For this purpose, ComBat is adapted for neuroimaging harmonisation in a study \citep{Fortinetal2018}, allowing scanner- and site-related variation to be reduced while biological variation is preserved. The harmonisation step also needs to be estimated without information from held-out test data. Deep representation learning is further developed through ASD-DiagNet, which is proposed in a study \citep{Eslamietal2019}. Representation learning and classification are combined, and an accuracy of approximately 70.3\% is reported. The small difference in performance indicates that classification \reviewhl{is also potentially influenced} by the way functional connectivity is represented.

This role of connectivity representation is examined more directly in a controlled benchmarking study \citep{Dadietal2019}. Multiple connectome representations and predictive models are compared under the same setting, and tangent-space connectivity is identified as a strong representation. With an appropriate connectivity representation, comparatively simple classifiers are shown to remain competitive with more complex nonlinear models. Sparse representation learning is further explored through ASD-SAENet, which is proposed in a study \citep{AlmuqhimSaeed2021}. A sparse autoencoder is combined with a deep neural network, and an accuracy of approximately 70.8\% is reported. Across these deep-learning approaches, the increase from 70.0\% to 70.3\% and then 70.8\% remains relatively small. This pattern suggests that deeper representation learning alone \reviewhl{does not necessarily fully address} the challenges of heterogeneous multisite rs-fMRI data.


\subsection{Evidence from Systematic Reviews}

The broader performance of rs-fMRI machine-learning methods is examined through a systematic review and meta-analysis in \cite{Santanaetal2022}. Moderate pooled sensitivity, specificity, and discrimination are reported, while considerable methodological heterogeneity is identified. Differences in preprocessing, sample selection, validation strategy, feature extraction, and classifier design make direct comparison between studies difficult.

A later systematic review of deep-learning and rs-fMRI biomarkers \cite{Hudaetal2024} reports continued methodological development but identifies similar limitations related to generalisability, reproducibility, and inconsistent evaluation. These reviews provide an important analytical perspective: higher accuracy reported by an individual model cannot be interpreted independently of the dataset definition and experimental pipeline used to obtain it.

Therefore, later ASD classification studies increasingly attempt to improve the representation itself rather than only increasing classifier depth.

\subsection{Multi-Atlas and Dynamic Connectivity Models}

Dependence on a single brain parcellation is addressed through MADE-for-ASD \cite{Liuetal2024MADE}. A multi-atlas deep ensemble is introduced, and approximately 75.2\% accuracy is reported. Different atlases describe the same brain according to different regional boundaries, so their predictions \reviewhl{potentially contain} complementary information. Multi-atlas fusion therefore provides a possible way to reduce dependence on one predefined parcellation.

However, the improvement is introduced together with a more complex ensemble architecture. It therefore becomes difficult to determine whether performance is improved because of atlas diversity, deep-model capacity, demographic information, or their interaction. This motivates controlled evaluations in which atlas fusion is varied while the remaining pipeline is kept unchanged.

Dynamic functional connectivity is later revisited through a GNN-LSTM framework \cite{Tangetal2025GNNLSTM}. Graph neural networks are combined with temporal sequence modelling, and approximately 80.4\% accuracy is reported on ABIDE-I. Compared with static approaches, connectivity variation across the scan is explicitly represented rather than being averaged into one matrix.

This development addresses an important limitation of static functional connectivity, but it introduces several additional decisions. Window length, overlap, temporal ordering, graph construction, and sequence-model complexity \reviewhl{all affect performance}. Therefore, the contribution of temporal variation \reviewhl{requires separation} from the contribution of the classifier through controlled ablation.

\subsection{Explainable and Multimodal ASD Models}

Interpretability receives greater attention in recent ASD classification research. Explainable artificial intelligence is applied to ABIDE fMRI data in a study \cite{Vidyaetal2025}, where highly predictive performance is reported under the stated filtering condition and influential brain regions are identified through multiple explanation methods. Such work demonstrates the potential of connecting prediction with biologically interpretable patterns.

However, unusually high performance on heterogeneous multisite data also increases the importance of strict validation. Feature selection, filtering, preprocessing, and cross-validation \reviewhl{are required to remain} independent of held-out data. The broader lesson from earlier benchmarking \cite{Dadietal2019} and meta-analysis \cite{Santanaetal2022} therefore remains relevant: reported performance \reviewhl{requires interpretation} together with the complete experimental procedure.

The progression toward richer representations continues through CAS-GNN \cite{Tangetal2026CASGNN}. Structural and functional MRI are integrated using cross-attention and subject-adaptive graph construction, and approximately 79.3\% accuracy is reported. Compared with earlier functional-connectivity-only approaches, multimodal information is used to represent each participant more comprehensively.

Although multimodal graph learning increases representational capacity, additional learned components also make generalisation more difficult to assess. Performance gains \reviewhl{therefore require evaluation} against cross-site stability, model complexity, and potential overfitting.

\subsection{Research Gap}

The literature shows two connected methodological developments. The first progresses from stable covariance estimation \cite{LedoitWolf2004}, geometric connectivity modelling \cite{Varoquauxetal2010}, and network-level representation \cite{Yeoetal2011} toward dynamic connectivity \cite{Hutchisonetal2013,LeonardiVanDeVille2015,Pretietal2017}, multisite analysis \cite{DiMartinoetal2014}, reproducibility \cite{Abrahametal2017}, deep learning \cite{Heinsfeldetal2018,Eslamietal2019,AlmuqhimSaeed2021}, multi-atlas fusion \cite{Liuetal2024MADE}, dynamic graph modelling \cite{Tangetal2025GNNLSTM}, and multimodal graph learning \cite{Tangetal2026CASGNN}. Each stage addresses a limitation of the previous representation, but greater representational complexity also introduces additional sources of variation and overfitting.

The second development concerns prediction reliability. Statistical comparison and confidence estimation are formalised in \cite{ClopperPearson1934,Demsar2006}, while harmonisation \cite{Johnsonetal2007,Fortinetal2018}, calibration \cite{Guoetal2017}, selective prediction \cite{GeifmanElYaniv2017}, explainability \cite{LundbergLee2017}, and decision-oriented evaluation \cite{VanCalsteretal2018} show that classification accuracy alone is insufficient for a screening-oriented system. Reliable probabilities, uncertainty-aware decisions, interpretable features, and performance on unseen acquisition sites are also required.

Despite these advances, these directions are commonly investigated separately. Dynamic connectivity, multi-atlas information, harmonisation, interpretability, calibration, selective prediction, and cross-site generalisation are rarely evaluated within the same controlled pipeline. It also remains unclear whether increasingly complex classifiers provide greater benefit than carefully designed functional-connectivity representations when identical validation conditions are used.

The present work is therefore designed around a controlled and reliability-oriented evaluation. Complementary functional-connectivity representations are compared under the same data partitions, tangent-space quantities are estimated from training data only, and cross-site performance is evaluated through leave-one-site-out validation. Structured low-rank connectivity patterns are used for interpretation, while calibrated probabilities and selective prediction are used to manage uncertainty. Each component is evaluated independently so that improvements \reviewhl{are attributable} to the corresponding methodological choice rather than to model complexity alone.
